\documentclass[11pt,a4paper]{article}

\usepackage[utf8]{inputenc}
\usepackage[T1]{fontenc}
\usepackage[margin=2.5cm]{geometry}
\usepackage{hyperref}
\usepackage{xcolor}
\usepackage{listings}
\usepackage{enumitem}
\usepackage{booktabs}
\usepackage{longtable}
\usepackage{amsmath,amssymb}
\usepackage{fancyhdr}
\usepackage{titlesec}

\definecolor{isls-blue}{HTML}{1a5276}
\definecolor{isls-orange}{HTML}{e67e22}
\definecolor{isls-green}{HTML}{27ae60}
\definecolor{isls-red}{HTML}{c0392b}
\definecolor{isls-gray}{HTML}{7f8c8d}
\definecolor{code-bg}{HTML}{f7f9fb}

\lstdefinestyle{rust}{
  language=C,
  basicstyle=\ttfamily\small,
  keywordstyle=\color{isls-blue}\bfseries,
  commentstyle=\color{isls-gray}\itshape,
  stringstyle=\color{isls-green},
  backgroundcolor=\color{code-bg},
  frame=single,
  rulecolor=\color{isls-gray!30},
  breaklines=true,
  showstringspaces=false,
  tabsize=4,
  morekeywords={fn,let,mut,pub,struct,impl,enum,use,mod,self,Self,
    crate,super,where,trait,for,in,if,else,match,return,async,await,
    Ok,Err,Some,None,Vec,String,HashMap,BTreeMap,BTreeSet,
    Result,Option,PathBuf,
    f64,f32,usize,bool,u32,u64,i64,Box,dyn,Arc},
}
\lstdefinestyle{bash}{
  language=bash,
  basicstyle=\ttfamily\small,
  backgroundcolor=\color{code-bg},
  frame=single,
  rulecolor=\color{isls-gray!30},
  breaklines=true,
  showstringspaces=false,
}
\lstdefinestyle{json}{
  basicstyle=\ttfamily\small,
  backgroundcolor=\color{code-bg},
  frame=single,
  rulecolor=\color{isls-gray!30},
  breaklines=true,
  showstringspaces=false,
}

\pagestyle{fancy}
\fancyhf{}
\fancyhead[L]{\small\textcolor{isls-gray}{ISLS I3 Spec --- Metamorphose-Kern}}
\fancyhead[R]{\small\textcolor{isls-gray}{v1.0.0 --- \today}}
\fancyfoot[C]{\thepage}

\titleformat{\section}
  {\Large\bfseries\color{isls-blue}}{\thesection}{1em}{}
\titleformat{\subsection}
  {\large\bfseries\color{isls-blue!80}}{\thesubsection}{1em}{}
\titleformat{\subsubsection}
  {\normalsize\bfseries\color{isls-blue!60}}{\thesubsubsection}{1em}{}

\hypersetup{
  colorlinks=true,
  linkcolor=isls-blue,
  urlcolor=isls-orange,
}

\begin{document}

\begin{titlepage}
\centering
\vspace*{2.5cm}

{\Huge\bfseries\color{isls-blue} ISLS I3 Specification}\\[0.5cm]
{\LARGE\color{isls-orange} Metamorphose-Kern}\\[0.3cm]
{\large\color{isls-gray} Constraint Spectroscopy, Norm-Injektion,\\
Solve-Coagula-Orchestrierung}\\[2cm]

{\large
\begin{tabular}{rl}
\textbf{System:}       & ISLS --- Intelligent Semantic Ledger Substrate \\
\textbf{Version:}      & v6.2 (I3) \\
\textbf{Author:}       & Sebastian Klemm \\
\textbf{Date:}         & \today \\
\textbf{Status:}       & Specification \\
\textbf{Depends:}      & I2 (complete) \\
\textbf{Branch:}       & \texttt{claude/implement-isls-i3-spec} \\
\textbf{Basis:}        & ISLS Metamorphose-Architektur (Formalisierung) \\
\end{tabular}
}

\vfill

{\small\color{isls-gray}
Mithras Substructure University\\
Repository: \texttt{https://github.com/LashSesh/graphity}\\
License: MIT
}
\end{titlepage}

\tableofcontents
\newpage


% ══════════════════════════════════════════════════════════════════════════════
\section{Executive Summary}
% ══════════════════════════════════════════════════════════════════════════════

After I2, ISLS generates, observes, and learns. But the 
learning is passive: the \"Olbohrinsel scrapes random keywords 
and hopes for useful norms. The system has no concept of 
\emph{what it needs}. It cannot analyze a target system and 
say ``I need norms for EventBus, IndicatorPipeline, and 
RiskManagement --- let me go find them.''

I3 installs three capabilities:

\begin{enumerate}[nosep]
\item \textbf{Constraint Spectroscopy:} Analyze a target system 
      (or a description), identify exactly which norm-classes 
      are missing, generate targeted scrape-keywords, and 
      fill the gaps chirurgically.
\item \textbf{Norm-Injektion:} Manually define norms as JSON 
      blueprints. The fitness system validates them 
      automatically --- bad norms die, good norms thrive.
\item \textbf{\texttt{isls evolve}:} A single CLI command that 
      orchestrates the full Solve-Coagula cycle: scrape an 
      existing system $\to$ identify gaps $\to$ optionally 
      fill gaps $\to$ re-forge with expanded norms. One 
      command to metamorphose.
\end{enumerate}

These three capabilities transform ISLS from a passive 
generator into an active, self-directed synthesis engine.


% ══════════════════════════════════════════════════════════════════════════════
\section{W1: Constraint Spectroscopy}
\label{sec:w1}
% ══════════════════════════════════════════════════════════════════════════════

\subsection{Problem}

The ``was fehlt mir?'' endpoint (S1c) checks against a 
hardcoded list of 20 common patterns. This is static and 
superficial. Constraint Spectroscopy replaces it with a 
dynamic, target-aware analysis.

\subsection{The Spectroscopy Pipeline}

\begin{lstlisting}[style=rust]
/// Analyze a target system and identify missing norm coverage.
pub struct SpectroscopyResult {
    /// Resonite classes found in the target.
    pub target_spectrum: Vec<ResoniteClass>,
    /// Resonite classes covered by existing norms.
    pub covered: Vec<ResoniteClass>,
    /// Gap: in target but not in norms.
    pub gaps: Vec<NormGap>,
    /// Suggested scrape keywords per gap.
    pub suggestions: Vec<ScrapeSuggestion>,
    /// Coverage ratio: |covered| / |target_spectrum|.
    pub coverage: f64,
}

pub struct NormGap {
    pub class: ResoniteClass,
    pub resonite_count: usize,  // how many resonites in target
    pub layers_affected: Vec<String>,
    pub priority: f64,  // count * layer_depth
}

pub struct ScrapeSuggestion {
    pub gap: String,
    pub keywords: Vec<String>,
    pub estimated_repos: usize,
}
\end{lstlisting}

\subsection{Resonite Classification}

\begin{lstlisting}[style=rust]
/// High-level classification of resonites into 
/// architectural pattern categories.
#[derive(Debug, Clone, Serialize, Deserialize, 
         Hash, Eq, PartialEq)]
pub enum ResoniteClass {
    CrudEntity,
    Authentication,
    Authorization,
    EventBus,
    MessageQueue,
    ConnectionPool,
    Caching,
    RateLimiting,
    CircuitBreaker,
    RetryPattern,
    Pagination,
    Search,
    FileUpload,
    ExportImport,
    Scheduling,
    Notification,
    Logging,
    Metrics,
    HealthCheck,
    Configuration,
    Migration,
    Testing,
    Docker,
    StateMachine,
    Workflow,
    IndicatorPipeline,
    SignalProcessing,
    RiskManagement,
    OrderExecution,
    DataVisualization,
    RealtimeWebSocket,
    GraphQLApi,
    GrpcService,
    CliInterface,
    PluginSystem,
    Custom(String),
}

/// Classify a resonite into a high-level class.
pub fn classify_resonite(r: &Resonite) -> Option<ResoniteClass> {
    match r {
        Resonite::Fn { name, .. } => {
            if name.starts_with("get_") || name.starts_with("list_") 
               || name.starts_with("create_") || name.starts_with("update_")
               || name.starts_with("delete_") {
                Some(ResoniteClass::CrudEntity)
            } else if name.contains("auth") || name.contains("login") 
                      || name.contains("token") {
                Some(ResoniteClass::Authentication)
            } else if name.contains("emit") || name.contains("subscribe") 
                      || name.contains("dispatch") {
                Some(ResoniteClass::EventBus)
            } else if name.contains("cache") || name.contains("invalidate") {
                Some(ResoniteClass::Caching)
            } else if name.contains("connect") || name.contains("pool") {
                Some(ResoniteClass::ConnectionPool)
            }
            // ... weitere Klassifikationen
            else { None }
        },
        Resonite::Type { name, .. } => {
            if name.contains("Event") || name.contains("Handler") {
                Some(ResoniteClass::EventBus)
            } else if name.contains("Cache") {
                Some(ResoniteClass::Caching)
            }
            // ... weitere
            else { None }
        },
        _ => None,
    }
}
\end{lstlisting}

\subsection{Keyword Generation}

\begin{lstlisting}[style=rust]
/// Generate scrape keywords for a gap.
pub fn suggest_keywords(gap: &NormGap) -> Vec<String> {
    let base = match &gap.class {
        ResoniteClass::EventBus => vec![
            "rust event bus async",
            "rust event driven architecture",
            "rust publish subscribe pattern",
        ],
        ResoniteClass::Caching => vec![
            "rust cache lru",
            "rust caching redis",
            "rust memoization",
        ],
        ResoniteClass::ConnectionPool => vec![
            "rust connection pool database",
            "rust deadpool bb8 r2d2",
        ],
        ResoniteClass::RateLimiting => vec![
            "rust rate limit tower",
            "rust rate limiter",
        ],
        ResoniteClass::IndicatorPipeline => vec![
            "rust technical indicators trading",
            "rust financial analysis ta",
        ],
        ResoniteClass::RiskManagement => vec![
            "rust risk management kelly criterion",
            "rust portfolio risk var",
        ],
        // ... 30+ Klassen mit je 2-3 Keywords
        _ => vec![],
    };
    base.into_iter().map(String::from).collect()
}
\end{lstlisting}

\subsection{CLI: \texttt{isls spectroscopy}}

\begin{lstlisting}[style=bash]
# Analyze a local project:
isls spectroscopy --path ./my-trading-system

ISLS Constraint Spectroscopy
---
Target: ./my-trading-system
Resonites: 234 (42 Fn, 89 Type, 67 Import, 12 Layer, 24 Relation)
Spectrum: 8 classes identified

Coverage: 3/8 (37.5%)
  [OK] CrudEntity       (covered by ISLS-NORM-0042)
  [OK] Authentication   (covered by ISLS-NORM-0088)
  [OK] Pagination       (covered by ISLS-NORM-0096)
  [GAP] EventBus        (14 resonites, 3 layers)
  [GAP] IndicatorPipeline (8 resonites, 2 layers)
  [GAP] RiskManagement  (6 resonites, 2 layers)
  [GAP] ConnectionPool  (4 resonites, 1 layer)
  [GAP] StateMachine    (3 resonites, 1 layer)

Suggested scrape campaigns:
  1. "rust event bus async" (priority: 42)
  2. "rust technical indicators" (priority: 16)
  3. "rust risk management kelly" (priority: 12)
  4. "rust connection pool deadpool" (priority: 4)
  5. "rust state machine" (priority: 3)

Run: isls spectroscopy --path ./my-trading-system --scrape
to automatically fill gaps.

# Analyze + auto-scrape:
isls spectroscopy --path ./my-trading-system --scrape

Filling 5 gaps...
  [1/5] EventBus: scraping 3 keywords, 15 repos...
        -> 3 new candidates, 1 promoted
  [2/5] IndicatorPipeline: scraping 2 keywords, 10 repos...
        -> 2 new candidates
  ...
Coverage after: 7/8 (87.5%)
  [GAP] StateMachine (still missing, only 2 repos found)
\end{lstlisting}

\subsection{API Endpoint}

\begin{lstlisting}[style=bash]
POST /api/discover/spectroscopy
Body: { "path": "/path/to/project" }
  or: { "url": "https://github.com/user/repo" }
  or: { "description": "Trading system with events and risk" }
Response: SpectroscopyResult as JSON

POST /api/discover/spectroscopy/fill
Body: { "gaps": ["EventBus", "Caching"], "repos_per_keyword": 5 }
Response: { "job_id": "...", "gaps": 2, "keywords": 5 }
(Progress via WebSocket, like mass-scrape)
\end{lstlisting}

\subsection{New Files}

\begin{longtable}{lp{10cm}}
\toprule
\textbf{File} & \textbf{Purpose} \\
\midrule
\texttt{isls-norms/src/spectroscopy.rs}
  & \texttt{ResoniteClass} enum, \texttt{classify\_resonite()}, 
    \texttt{spectroscopy()}, \texttt{suggest\_keywords()}, 
    \texttt{SpectroscopyResult}. \\
\bottomrule
\end{longtable}

\subsection{Modified Files}

\begin{longtable}{lp{10cm}}
\toprule
\textbf{File} & \textbf{Change} \\
\midrule
\texttt{isls-norms/src/lib.rs}
  & \texttt{pub mod spectroscopy;} \\
\texttt{isls-cli/src/main.rs}
  & Add \texttt{isls spectroscopy} subcommand with 
    \texttt{--path}, \texttt{--url}, \texttt{--scrape} flags. \\
\texttt{isls-gateway/src/discover.rs}
  & Replace hardcoded \texttt{gaps()} with 
    \texttt{spectroscopy()} call. Add 
    \texttt{/api/discover/spectroscopy} and 
    \texttt{/api/discover/spectroscopy/fill} endpoints. \\
\texttt{isls-gateway/src/lib.rs}
  & Route new endpoints. \\
\bottomrule
\end{longtable}

\subsection{W1 Verification}

\begin{enumerate}[nosep]
\item \texttt{isls spectroscopy --path .} on the ISLS repo 
      itself produces a spectrum.
\item Gaps are identified (ISLS has no EventBus, no Caching).
\item Keywords are suggested for each gap.
\item \texttt{--scrape} flag triggers targeted scraping.
\item Studio: ``was fehlt mir?'' now uses spectroscopy instead 
      of hardcoded list.
\end{enumerate}


% ══════════════════════════════════════════════════════════════════════════════
\section{W2: Norm-Injektion}
\label{sec:w2}
% ══════════════════════════════════════════════════════════════════════════════

\subsection{Problem}

Norms can only enter the system through auto-discovery 
(observation $\to$ candidate $\to$ promotion). There is no 
way for the architect to directly define a norm. This blocks 
two use cases:
\begin{itemize}[nosep]
\item Novel patterns that don't exist in any scraped repo.
\item Domain-specific patterns from proprietary knowledge.
\end{itemize}

\subsection{Norm Blueprint Format}

\begin{lstlisting}[style=json]
{
  "id": "ISLS-NORM-INJECT-001",
  "name": "EventDrivenArchitecture",
  "level": "System",
  "description": "Event bus with handlers, dispatcher, store",
  "layers": [
    "EventDefinition",
    "EventBus",
    "EventHandler",
    "EventDispatcher",
    "EventStore"
  ],
  "expected_resonites": [
    { "kind": "Fn", "pattern": "emit_*", "arity": 2 },
    { "kind": "Fn", "pattern": "subscribe_*", "arity": 2 },
    { "kind": "Fn", "pattern": "dispatch", "arity": 1 },
    { "kind": "Fn", "pattern": "replay_events", "arity": 1 },
    { "kind": "Type", "pattern": "Event", "type_kind": "Enum" },
    { "kind": "Type", "pattern": "*Handler", "type_kind": "Trait" },
    { "kind": "Type", "pattern": "EventStore", "type_kind": "Struct" }
  ],
  "activation_keywords": [
    "event", "event-driven", "publish", "subscribe",
    "dispatch", "handler", "event sourcing"
  ],
  "constraints": {
    "min_layers": 3,
    "requires": ["ISLS-NORM-INFRA-DB"]
  }
}
\end{lstlisting}

\subsection{Implementation}

\begin{lstlisting}[style=rust]
/// A manually injected norm blueprint.
#[derive(Debug, Clone, Serialize, Deserialize)]
pub struct NormBlueprint {
    pub id: String,
    pub name: String,
    pub level: String,           // Atom, Molecule, System
    pub description: String,
    pub layers: Vec<String>,
    pub expected_resonites: Vec<ResonitePattern>,
    pub activation_keywords: Vec<String>,
    pub constraints: BlueprintConstraints,
}

#[derive(Debug, Clone, Serialize, Deserialize)]
pub struct ResonitePattern {
    pub kind: String,       // "Fn", "Type", "Import"
    pub pattern: String,    // glob pattern, e.g. "emit_*"
    pub arity: Option<usize>,
    pub type_kind: Option<String>,
}

#[derive(Debug, Clone, Serialize, Deserialize)]
pub struct BlueprintConstraints {
    pub min_layers: Option<usize>,
    pub requires: Vec<String>,    // norm IDs that must co-activate
}

/// Inject a norm blueprint into the registry.
pub fn inject_norm(
    registry: &mut NormRegistry,
    blueprint: NormBlueprint,
) -> Result<String, String> {
    // 1. Validate: ID must start with "ISLS-NORM-INJECT-"
    if !blueprint.id.starts_with("ISLS-NORM-INJECT-") {
        return Err("Injected norms must have ID prefix \
                    ISLS-NORM-INJECT-".to_string());
    }
    
    // 2. Check for duplicates
    if registry.get_norm(&blueprint.id).is_some() {
        return Err(format!("Norm {} already exists", blueprint.id));
    }
    
    // 3. Convert to internal Norm format
    let norm = Norm {
        id: blueprint.id.clone(),
        name: blueprint.name,
        level: parse_level(&blueprint.level),
        layers: blueprint.layers.iter()
            .map(|l| parse_layer(l)).collect(),
        builtin: false,
        source: NormSource::Injected,
        activation_keywords: blueprint.activation_keywords,
        // ... map remaining fields
    };
    
    // 4. Add to registry with default fitness 0.5
    registry.add_norm(norm);
    
    // 5. Initialize fitness
    fitness_store.set_fitness(&blueprint.id, 0.5);
    
    Ok(blueprint.id)
}
\end{lstlisting}

\subsection{CLI}

\begin{lstlisting}[style=bash]
# Inject a single norm:
isls norms inject ./event_driven.json

[OK] Injected ISLS-NORM-INJECT-001 "EventDrivenArchitecture"
     Fitness: 0.50 (neutral, will be validated by usage)
     Keywords: event, event-driven, publish, subscribe, ...

# List injected norms:
isls norms list --source injected

ISLS-NORM-INJECT-001  EventDrivenArchitecture  injected  System
ISLS-NORM-INJECT-002  IndicatorPipeline        injected  Organism

# Remove an injected norm:
isls norms remove ISLS-NORM-INJECT-001

[OK] Removed ISLS-NORM-INJECT-001. Fitness data preserved.
\end{lstlisting}

\subsection{Self-Correction}

Injected norms are subject to the same fitness tracking as 
auto-discovered norms (I1/I2). After each generation:

\begin{itemize}[nosep]
\item Compile success: $\phi \to \phi \cdot 0.9 + r \cdot 0.1$
\item $r = \text{codematrix\_resonance}$ if compiled, 
      $r = 0$ if failed.
\item After 20 failed generations: $\phi < 0.06$, 
      effectively dead.
\item No special treatment. Injected norms must earn their 
      fitness like any other norm.
\end{itemize}

\subsection{New Files}

\begin{longtable}{lp{10cm}}
\toprule
\textbf{File} & \textbf{Purpose} \\
\midrule
\texttt{isls-norms/src/injection.rs}
  & \texttt{NormBlueprint}, \texttt{inject\_norm()}, 
    \texttt{remove\_injected()}, blueprint validation. \\
\bottomrule
\end{longtable}

\subsection{Modified Files}

\begin{longtable}{lp{10cm}}
\toprule
\textbf{File} & \textbf{Change} \\
\midrule
\texttt{isls-norms/src/lib.rs}
  & \texttt{pub mod injection;} Add 
    \texttt{NormSource::Injected} variant. \\
\texttt{isls-cli/src/main.rs}
  & \texttt{isls norms inject <file>}, 
    \texttt{isls norms remove <id>}, 
    \texttt{isls norms list --source} filter. \\
\bottomrule
\end{longtable}

\subsection{W2 Verification}

\begin{enumerate}[nosep]
\item Create a valid blueprint JSON.
\item \texttt{isls norms inject blueprint.json} succeeds.
\item \texttt{isls norms list} shows the injected norm.
\item A forge with matching keywords activates the injected norm.
\item An invalid blueprint (bad JSON, missing fields) is rejected.
\item After 5 failed generations: fitness drops below 0.3.
\end{enumerate}


% ══════════════════════════════════════════════════════════════════════════════
\section{W3: \texttt{isls evolve} --- Solve-Coagula Orchestrierung}
\label{sec:w3}
% ══════════════════════════════════════════════════════════════════════════════

\subsection{Problem}

The Solve-Coagula cycle currently requires manual steps:
\texttt{isls scrape} $\to$ (think about gaps) $\to$ 
\texttt{isls forge-chat}. There is no single command 
that orchestrates the full metamorphose.

\subsection{The \texttt{evolve} Command}

\begin{lstlisting}[style=bash]
isls evolve \
    --from ./current-version \
    --delta "Add real-time WebSocket notifications and 
             event-driven architecture" \
    --output ./next-version \
    --scrape-gaps \
    --ollama
\end{lstlisting}

This single command:

\begin{enumerate}[nosep]
\item \textbf{Solve:} Scrapes \texttt{./current-version}, 
      extracts topology, feeds to norm system.
\item \textbf{Spectroscopy:} Analyzes the delta description 
      + existing topology. Identifies missing norms.
\item \textbf{Fill (optional, \texttt{--scrape-gaps}):} 
      Runs chirurgical scraping for each gap.
\item \textbf{Merge:} Combines the delta description with 
      the existing topology to form a complete generation 
      description.
\item \textbf{Coagula:} Runs \texttt{forge-chat} with the 
      merged description and the enriched norm library.
\item \textbf{Report:} Shows before/after comparison.
\end{enumerate}

\subsection{Implementation}

\begin{lstlisting}[style=rust]
pub async fn cmd_evolve(
    from: PathBuf,
    delta: String,
    output: PathBuf,
    scrape_gaps: bool,
    oracle: Box<dyn Oracle>,
) -> Result<EvolveReport, Box<dyn Error>> {
    // 1. SOLVE: Extract topology of existing system
    tracing::info!("Solve: scraping {}", from.display());
    let topo = scrape_path(&from)?;
    let resonites = extract_resonites_from_topology(&topo);
    observe_and_learn(&topo, &from)?;
    
    // 2. SPECTROSCOPY: Identify gaps
    tracing::info!("Spectroscopy: analyzing gaps");
    let spectrum = spectroscopy(&resonites, &load_norms()?);
    tracing::info!("Coverage: {:.0}% ({} gaps)", 
        spectrum.coverage * 100.0, spectrum.gaps.len());
    
    // 3. FILL GAPS (optional)
    if scrape_gaps && !spectrum.gaps.is_empty() {
        tracing::info!("Filling {} gaps via targeted scraping", 
            spectrum.gaps.len());
        for suggestion in &spectrum.suggestions {
            for keyword in &suggestion.keywords {
                // GitHub search + scrape top 5
                mass_scrape_keyword(keyword, 5).await?;
            }
        }
        // Re-evaluate coverage
        let new_spectrum = spectroscopy(&resonites, &load_norms()?);
        tracing::info!("Coverage after fill: {:.0}%", 
            new_spectrum.coverage * 100.0);
    }
    
    // 4. MERGE: Build combined description
    let existing_entities = extract_entities_from_topology(&topo);
    let merged_description = format!(
        "{}\n\nExisting system has: {}. \
         Retain all existing functionality.",
        delta,
        existing_entities.join(", ")
    );
    
    // 5. COAGULA: Re-forge
    tracing::info!("Coagula: forging next generation");
    let result = forge_chat(
        &merged_description, 
        &output, 
        oracle
    ).await?;
    
    // 6. REPORT
    let report = EvolveReport {
        from: from.clone(),
        output: output.clone(),
        delta: delta.clone(),
        solve_resonites: resonites.len(),
        gaps_before: spectrum.gaps.len(),
        gaps_after: if scrape_gaps { 
            spectroscopy(&resonites, &load_norms()?)
                .gaps.len() 
        } else { spectrum.gaps.len() },
        coverage_before: spectrum.coverage,
        forge_result: result,
    };
    
    print_evolve_report(&report);
    Ok(report)
}
\end{lstlisting}

\subsection{Report Output}

\begin{lstlisting}[style=bash]
isls evolve --from ./trading-v1 \
    --delta "Add WebSocket notifications" \
    --output ./trading-v2 --scrape-gaps --ollama

ISLS Evolve: Solve-Coagula Cycle
================================

[Solve] Scraped ./trading-v1
  234 resonites, 8 classes, 42 functions, 89 types

[Spectroscopy] Coverage: 37.5% (5 gaps)
  [GAP] EventBus (priority: 42)
  [GAP] RealTimeWebSocket (priority: 28)
  [GAP] ConnectionPool (priority: 4)
  [GAP] IndicatorPipeline (priority: 16)
  [GAP] RiskManagement (priority: 12)

[Fill] Targeted scraping for 5 gaps...
  EventBus: 15 repos -> 3 candidates, 1 norm
  WebSocket: 10 repos -> 2 candidates
  ConnectionPool: 10 repos -> 1 candidate, 1 norm
  IndicatorPipeline: 10 repos -> 2 candidates
  RiskManagement: 5 repos -> 1 candidate
  Coverage after: 62.5% (3 remaining gaps)

[Coagula] Forging next generation...
  Description: "Add WebSocket notifications. Existing 
  system has: Trade, Portfolio, Signal, Order, Indicator.
  Retain all existing functionality."
  
  35 files, 1580 LOC, 7 LLM calls, 0 coagula, 189s
  cargo check: PASSED

[Report]
  Before: ./trading-v1 (234 resonites, 37.5% coverage)
  After:  ./trading-v2 (298 resonites, 62.5% coverage)
  Delta:  +64 resonites, +25% coverage, +WebSocket
  Norms used: 38 (33 existing + 2 new auto + 3 injected)
\end{lstlisting}

\subsection{Studio Integration}

In the Studio Erschaffen mode, the user can type:

\begin{lstlisting}[style=bash]
evolve ./trading-v1 with WebSocket notifications
\end{lstlisting}

The Studio detects the ``evolve'' keyword, extracts the 
path and delta, and runs the evolve pipeline. Progress 
appears in the chat stream (Solve $\to$ Spectroscopy $\to$ 
Fill $\to$ Coagula $\to$ Report).

\subsection{Modified Files}

\begin{longtable}{lp{10cm}}
\toprule
\textbf{File} & \textbf{Change} \\
\midrule
\texttt{isls-cli/src/main.rs}
  & Add \texttt{isls evolve} subcommand with 
    \texttt{--from}, \texttt{--delta}, \texttt{--output}, 
    \texttt{--scrape-gaps}, \texttt{--ollama} flags. \\
\texttt{isls-cli/src/cmd\_evolve.rs} (NEW)
  & \texttt{cmd\_evolve()} implementation. \\
\texttt{isls-gateway/src/lib.rs}
  & \texttt{POST /api/evolve} endpoint (for Studio). \\
\bottomrule
\end{longtable}

\subsection{W3 Verification}

\begin{enumerate}[nosep]
\item \texttt{isls evolve --from . --delta "add health checks" 
      --output ./evolved} runs without error.
\item Solve phase: scrapes and observes the source project.
\item Spectroscopy phase: identifies gaps and suggests keywords.
\item \texttt{--scrape-gaps}: triggers targeted scraping.
\item Coagula phase: produces a compilable project in 
      \texttt{--output}.
\item Report shows before/after coverage comparison.
\end{enumerate}


% ══════════════════════════════════════════════════════════════════════════════
\section{W4: Studio UX f\"ur Spectroscopy + Evolve}
\label{sec:w4}
% ══════════════════════════════════════════════════════════════════════════════

\subsection{Entdecken-Modus: ``was fehlt mir?'' Upgrade}

Replace the hardcoded gap response with spectroscopy output. 
When the user types ``was fehlt mir?'' in Entdecken mode:

\begin{enumerate}[nosep]
\item \texttt{GET /api/discover/spectroscopy} (no target: 
      analyzes norm library breadth against all known classes).
\item Render gaps as cards with [Scrape] buttons (existing UI).
\item Add coverage percentage at the top.
\end{enumerate}

\subsection{Entdecken-Modus: Target Spectroscopy}

When user types ``analysiere <path or url>'':

\begin{enumerate}[nosep]
\item \texttt{POST /api/discover/spectroscopy} with path/url.
\item Render full SpectroscopyResult: spectrum, covered, gaps, 
      suggestions.
\item Each gap has a [F\"ullen] button that triggers 
      \texttt{/api/discover/spectroscopy/fill}.
\end{enumerate}

\subsection{Erschaffen-Modus: Evolve}

When user types ``evolve <path> with <delta>'':

\begin{enumerate}[nosep]
\item Parse path and delta from message.
\item \texttt{POST /api/evolve} with path, delta, scrape\_gaps.
\item Render progress as process messages (Solve $\to$ 
      Spectroscopy $\to$ Fill $\to$ Coagula $\to$ Report).
\end{enumerate}

\subsection{Modified Files}

\begin{longtable}{lp{10cm}}
\toprule
\textbf{File} & \textbf{Change} \\
\midrule
\texttt{isls-gateway/src/static/studio.html}
  & Parse ``analysiere'' and ``evolve'' keywords. Render 
    spectroscopy results and evolve progress. \\
\bottomrule
\end{longtable}


% ══════════════════════════════════════════════════════════════════════════════
\section{Dependency Graph}
% ══════════════════════════════════════════════════════════════════════════════

\begin{verbatim}
W1 (Constraint Spectroscopy)
  |
  +---> W2 (Norm-Injektion)  [independent]
  |
  +---> W3 (isls evolve)     [uses W1 spectroscopy]
          |
          v
        W4 (Studio UX)       [uses W1 + W3 endpoints]
\end{verbatim}

W1 first (foundation). W2 independent (can be parallel). 
W3 depends on W1. W4 depends on W1 + W3.

Recommended: W1 $\to$ W2 $\to$ W3 $\to$ W4.


% ══════════════════════════════════════════════════════════════════════════════
\section{Constraints for Claude Code}
% ══════════════════════════════════════════════════════════════════════════════

\begin{enumerate}
\item \textbf{Spectroscopy} uses \texttt{isls\_reader} for parsing 
      and \texttt{isls\_norms} for norm registry access. No new 
      external dependencies.
\item \textbf{ResoniteClass} classification is keyword-based 
      (function/type name matching). Not LLM-based. Deterministic.
\item \textbf{Norm-Injektion} validates JSON strictly. Missing 
      required fields $\to$ reject. Invalid level $\to$ reject. 
      Duplicate ID $\to$ reject.
\item \textbf{Injected norms} get ID prefix 
      \texttt{ISLS-NORM-INJECT-}. Never overwrite builtin or 
      auto-discovered norms.
\item \textbf{\texttt{isls evolve}} is a composition of existing 
      functions: \texttt{scrape} + \texttt{spectroscopy} + 
      \texttt{mass\_scrape} + \texttt{forge\_chat}. No new 
      generation logic. Only orchestration.
\item \textbf{The ``was fehlt mir?''} endpoint in discover.rs 
      must be replaced with spectroscopy call, not kept as 
      parallel implementation.
\item \textbf{Do not modify} isls-merkaba, isls-reader, 
      isls-code-topo, isls-hypercube, isls-types, isls-chat, 
      isls-forge-llm.
\item All existing tests MUST pass.
\item Branch: \texttt{claude/implement-isls-i3-spec}.
\item Commit format: \texttt{I3/W\{n\}: <description>}.
\end{enumerate}


% ══════════════════════════════════════════════════════════════════════════════
\section{Acceptance Criteria}
% ══════════════════════════════════════════════════════════════════════════════

\begin{longtable}{clp{9cm}}
\toprule
\textbf{\#} & \textbf{Pass} & \textbf{Criterion} \\
\midrule
1 & & \texttt{cargo build --workspace} compiles. \\
2 & & \texttt{cargo test --workspace} passes. \\
3 & & \texttt{isls spectroscopy --path .} produces coverage report. \\
4 & & \texttt{isls norms inject blueprint.json} adds norm to registry. \\
5 & & \texttt{isls norms list --source injected} shows injected norms. \\
6 & & \texttt{isls evolve --from . --delta "add X" --output ./out} 
      runs full cycle. \\
7 & & Studio ``was fehlt mir?'' uses spectroscopy. \\
8 & & Studio ``evolve'' triggers full cycle with progress. \\
9 & & Invalid blueprint JSON is rejected with clear error. \\
10 & & All existing tests pass (backward compat). \\
\bottomrule
\end{longtable}

I3 is passed when criteria 1--6 and 10 are fulfilled. 
Criteria 7--8 require Studio testing.


\vfill
\begin{center}
\color{isls-gray}\rule{0.5\textwidth}{0.4pt}\\[0.5em]
{\small End of I3 Specification.\\[0.3em]
Das System wei{\ss} jetzt was es nicht wei{\ss}.\\
Und es wei{\ss} wo es suchen muss.\\
Und es kann sich selbst neu drucken.}
\end{center}

\end{document}
